\scanpage{24}
% archive sheet number: 23
qui contient $\unsure{\widehat{\mathscr E}|\widehat{\mathcal E}}$ comme sous-extension. On désigne par $\unsure{\mathbb\Gamma|\Pi}$ le commutant de $\unsure{\mathfrak S_3|\mathscr S_3|G_3}$ dans $\widehat{\widehat\Gamma}$ \inserted{(contenant $\unsure{\dot\tau|\tau}$)}, et par $\unsure{\mathbb E|E}$ l'extension correspondante de $\mathbb\Gamma$ par $\widehat\pi$. Notons déjà que
\[
\left\{\begin{aligned}
\mathbb\Gamma\cap\Gamma_{\unsure{0,3}}&=\{1,\tau\}\\
\mathbb\Gamma\cap\mathfrak S_3&=\{1\}
\end{aligned}\right.
\]
\margin{\deleted{\unsure{NB}}}
% slanted
\margin{\uline{NB} $\widehat{\widehat\Gamma}\supset\mathfrak S_3\times\mathbb\Gamma$; $\mathfrak S_3\times\widehat{\widehat\Gamma}^{!}\supset\widehat{\widehat\Gamma}$. --- j'ignore si cette inclusion est une égalité, i.e. si $\mathbb\Gamma\hookrightarrow\widehat{\widehat\Gamma}^{!}$ est une égalité.}

On veut donner d'abord une description plus symétrique des éléments de $\widehat{\widehat{\mathscr E}}^{!}$, sans faire jouer un rôle particulier à $L_0$ engendré par $l_0$ $(L_0=\widehat k_0(\widehat{\mathbb Z}))$. Un $u\in\widehat{\widehat{\mathscr E}}$ est \inserted{peut être} défini par un multiplicateur
\[
\unsure{p}=\chi(u)\in\widehat{\mathbb Z}^{*}\tag*{(19)}
\]
et des
\[
g_i\in\widehat\pi\qquad i=0,1,\infty\tag*{(20)}
\]
pour
\[
u(l_i)=\operatorname{int}(g_i^{\unsure{-1|1}})l_i^p\tag*{(21)}
\]
les données (19), (20) étant liées par la relation
\[
(\operatorname{int}(g_\infty^{-1})l_\infty^p)(\operatorname{int}g_1^{-1}\cdot l_1^p)(\operatorname{int}(g_0^{-1})l_0^p)=1\tag*{(22)}
\]
et par le fait que les $u(l_i)$ définis par (21) \uline{engendrent} $\widehat\pi$ topologiquement. Pour $u$ fixé, $p$
